\section{Introduction}
\label{sec:introduction}

This sample paper is also the manual of the \labname{} arXiv class.
Each section lives in its own file under \texttt{sections/}, drawings under \texttt{figures/}, table bodies under \texttt{tables/}, the bibliography under \texttt{bib/}, and the macros a paper adds in \texttt{commands.tex}.
\texttt{AGENTS.md} lists the conventions, and this section shows the two that matter most in running text.

A citation that works as a noun, the subject or the object of the sentence, uses \verb|\citet|.
For example, \citet{vaswani2017attention} introduced the Transformer.
A citation that only supports a claim or a named thing uses \verb|~\citep|, placed right after what it supports.
Large language models~\citep{brown2020language} can reason step by step when prompted to~\citep{wei2022chain}.
Both commands keep working when the \texttt{numbers} class option switches the citations from author-year to numbers.

Cross-references use \verb|\cref|, which prints the name with the number, so the text never types ``Figure'' or a number by hand.
\Cref{sec:method} describes the method, \cref{fig:overview} gives its overview, and \cref{tab:main-results} reports the results.
